\subsection{The module associated to an endomorphism}

Let A be a commutative ring, let M be an A-module, and let u be an A-endomorphism of M. Recall (III, p.~538) that the map \((p(\mathbf{X}), x) \mapsto p(u) \cdot x\) from \(A[X] \times M\) into M makes M an A{[}X{]}-module, written \(M_{u}\) . Recall also (III, pp.~538 and 539) that if M{[}X{]} denotes the A{[}X{]}-module obtained from M by extension of scalars from A to A{[}X{]}, and if \(\bar{u}\) denotes the A{[}X{]}-endomorphism of M{[}X{]} induced by u, then there is an exact sequence of A{[}X{]}-modules \(^{1}\)

\[
0 \to \mathbf {M} [ \mathbf {X} ] \stackrel {\psi} {\to} \mathbf {M} [ \mathbf {X} ] \stackrel {\varphi} {\to} \mathbf {M} _ {u} \to 0,\tag{1}
\]

where \(\varphi(p(X) \otimes x) = p(u).x\) and \(\psi = X - \bar{u}\) .

An endomorphism \(u\) of an A-module M and an endomorphism \(u'\) of an A-module \(M'\) are said to be similar if there exists an isomorphism \(g\) from M onto \(M'\) such that \(u' \circ g = g \circ u\) , that is (III, p.~540, Prop. 19) an isomorphism \(g\) from \(M_u\) onto \(M_{u'}'\) . If M (resp. \(M'\) ) is free on the finite basis B (resp. \(B'\) ), and if \(M(u)\) (resp. \(M(u')\) ) is the matrix of \(u\) (resp. \(u'\) ) with respect to B (resp. \(B'\) ), then \(u\) and \(u'\) are similar if and only if \(M(u)\) and \(M(u')\) are similar matrices (II, p.~356, Def. 6). The characteristic polynomials (III, p.~541, Def. 3) of two similar endomorphisms of finitely generated free modules are equal (III, p.~540, Prop. 19).

Let K be a commutative field; then any pair \((E, u)\) consisting of a vector space E over K and an endomorphism u of E corresponds to a K{[}X{]}-module \(E_{u}\) . Since the ring K{[}X{]} is a principal ideal domain (IV, p.~11, Prop. 11), the results of the preceding sections can be applied to \(E_{u}\) .

Let us first show how to translate certain notions from the language of modules to that of endomorphisms of vector spaces :

« V is a submodule of \(E_{u}\) » means : « V is a vector subspace of E closed under u ».

« V is a cyclic submodule of \(E_{u}\) » means : « there exists \(x \in V\) such that the vector subspace V is spanned by the elements \(u^{i} \cdot x (i \in N)\) ». Then V is said to be cyclic (with respect to u) and x is called a generator.

« V is an indecomposable submodule of \(E_{u}\) » means : « V is nonzero and is not the direct sum of two nonzero subspaces each closed under u ».

« a is the annihilator of the submodule V » means : « a is the ideal consisting of those polynomials \(p(\mathbf{X}) \in \mathbf{K}[\mathbf{X}]\) such that \(p(u).x = 0\) for all \(x \in V\) ».

The monic polynomial g such that a is equal to the principal ideal \((g)\) is then called the minimal polynomial of the restriction of u to V.

« \(E_{u}\) is cyclic with annihilator \(a = (g)\) »

\[
\left(\text { with } g (\mathbf {X}) = \mathbf {X} ^ {n} + \alpha_ {n - 1} \mathbf {X} ^ {n - 1} + \dots + \alpha_ {0}\right)
\]

\(^{1}\) The injectivity of \(\psi\) , which is not stated in Prop. 18 of III, p.~539, is proved as follows: in the notation of loc. cit., we have

\[
\psi \left(\sum \left(\mathrm{X} ^ {k} \otimes x _ {k}\right)\right) = \sum \mathrm{X} ^ {k} \otimes \left(x _ {k - 1} - u (x _ {k})\right).
\]

If \(\sum X^k \otimes x_k\) belongs to the kernel of \(\psi\) , then it follows that \(x_{k-1} = u(x_k)\) for all \(k\) , and the \(x_k\) are all zero, since the family \((x_k)\) has finite support.

means: « there exists \(x \in E\) such that \((u^{i} \cdot x)\) ( \(0 \leqslant i \leqslant n - 1\) ) is a basis of the vector space E, and \(g(u) \cdot x = 0\) ». In other words, we can find a basis of E such that the matrix U of u with respect to this basis is

\[
U = \left( \begin{array}{c c c c} 0 & 0 & 0 \dots 0 & - \alpha_ {0} \\ 1 & 0 & 0 \dots 0 & - \alpha_ {1} \\ 0 & 1 & 0 \dots 0 & - \alpha_ {2} \\ \cdots & \cdots & \cdots & \cdots \\ 0 & 0 & 0 \dots 0 & - \alpha_ {n - 2} \\ 0 & 0 & 0 \dots 1 & - \alpha_ {n - 1} \end{array} \right).\tag{2}
\]

« \(E_{u}\) is a torsion module » means, by the characterisation of cyclic torsion modules given above : « every cyclic submodule of \(E_{u}\) is finite dimensional over K ». In particular :

« \(E_{u}\) is a finitely generated torsion module » means : « E is finite dimensional over K ».

